\documentclass[11pt,a4paper]{article}

% --------------------------------------------------
% PAQUETES
% --------------------------------------------------
\usepackage[spanish]{babel}
\usepackage[utf8]{inputenc}
\usepackage[T1]{fontenc}
\usepackage{lmodern}

\usepackage[a4paper,
            left=2.8cm,
            right=2.8cm,
            top=3cm,
            bottom=3cm]{geometry}

\usepackage{setspace}
\usepackage{titlesec}
\usepackage{enumitem}
\usepackage{xcolor}
\usepackage{tcolorbox}
\usepackage{graphicx}
\usepackage{hyperref}
\usepackage{fancyhdr}
\usepackage{booktabs}
\usepackage{parskip}

% --------------------------------------------------
% CONFIGURACIÓN
% --------------------------------------------------

\definecolor{mainblue}{RGB}{25,55,109}
\definecolor{lightblue}{RGB}{235,242,250}

\hypersetup{
    colorlinks=true,
    linkcolor=mainblue,
    urlcolor=mainblue
}

\setstretch{1.15}

\pagestyle{fancy}
\fancyhf{}
\fancyhead[L]{AP 08 - Desglose de un incidente complejo}
\fancyhead[R]{Xavier Margalef Riestra}
\fancyfoot[C]{\thepage}

\titleformat{\section}
{\Large\bfseries\color{mainblue}}
{\thesection.}{0.5em}{}

\titleformat{\subsection}
{\large\bfseries}
{\thesubsection.}{0.5em}{}

\titleformat{\subsubsection}
{\normalsize\bfseries}
{\thesubsubsection.}{0.5em}{}

\tcbset{
    colback=lightblue,
    colframe=mainblue,
    arc=2mm
}

% --------------------------------------------------
% DOCUMENTO
% --------------------------------------------------

\begin{document}

% ==================================================
% PORTADA
% ==================================================

\begin{titlepage}

\centering

\vspace*{2cm}

{\Huge\bfseries Desglose de un Incidente Complejo \par}

\vspace{0.5cm}

{\Large Actividad Práctica AP 08\par}

\vspace{2cm}

\rule{\textwidth}{1.2pt}

\vspace{1cm}

{\Large\bfseries Seguridad de la Información y Ciberseguridad}

\vspace{1cm}

\rule{\textwidth}{1.2pt}

\vfill

\begin{tabular}{ll}
\textbf{Alumno:} & Xavier Margalef Riestra \\
\\
\textbf{Actividad:} & AP 08 - Desglosar un incidente complejo \\
\\
\textbf{Fecha de entrega:} & 14 de Julio de 2026 \\
\end{tabular}

\vfill

{\large Informe Técnico}

\vspace{1cm}

\end{titlepage}

\tableofcontents
\newpage

% ==================================================
% CONTENIDO
% ==================================================

\section{Reacción Inmediata}

\subsection{¿Qué tipo de auditoría debe ejecutarse en este mismo momento y por qué?}

\textbf{Respuesta a Incidentes dentro de Seguridad de la Información y Ciberseguridad}

Habría que realizar un \textbf{DFIR (Digital Forensics and Incident Response)}, ya que entra en acción después de que se ha producido un incidente o ciberataque. En esta ocasión el ataque ya ha sido perpetrado y ha generado consecuencias, por lo que debemos identificar y recopilar información sobre lo ocurrido y determinar cómo reaccionar.

\begin{tcolorbox}
Objetivos principales del DFIR:
\end{tcolorbox}

\begin{itemize}
    \item Identificar el origen del ataque.
    \item Reconstruir la cadena de eventos (\textit{timeline}).
    \item Preservar evidencias digitales (logs, memoria y discos).
    \item Establecer el alcance real de la exfiltración de datos.
\end{itemize}

En este caso concreto, el \textbf{DFIR} es crítico porque existe:

\begin{itemize}
    \item Fuga masiva de datos.
    \item Impacto legal derivado de la investigación gubernamental.
    \item Riesgo de pérdida de licencias de Visa y Mastercard.
\end{itemize}

\subsubsection*{CAATs utilizados en esta fase}

\begin{itemize}
    \item \textbf{ELK Stack / Splunk} → análisis de grandes volúmenes de logs.
    \item \textbf{Wireshark} → análisis de tráfico de red sospechoso.
    \item \textbf{Autopsy} → análisis forense de discos.
    \item \textbf{Snort / OSSEC} → detección de intrusiones previas.
\end{itemize}

% ==================================================

\section{Análisis}

\textit{Si GlobalPay hubiera hecho bien su trabajo antes del ataque, ¿qué tres tipos específicos de auditorías preventivas habrían detectado y evitado los pasos 1, 2 y 3 del atacante? Relaciona cada auditoría con su paso correspondiente.}

\vspace{0.5em}

Aquí aplicamos conceptos de Gestión de Proyectos IT, SDLC, Seguridad de la Información y Ciberseguridad.

% --------------------------------------------------

\subsection{Paso 1: XSS en portal web obsoleto → robo de sesión}

Se utilizaría una auditoría de \textbf{Aplicaciones y Desarrollo (SDLC)} mediante análisis web enfocado en \textbf{Seguridad en Aplicaciones y APIs (OWASP)}.

\textbf{Habría detectado:}

\begin{itemize}
    \item Vulnerabilidades OWASP Top 10 (XSS).
    \item Falta de validación de entradas.
    \item Ausencia de pruebas de seguridad dentro del ciclo de vida del software.
\end{itemize}

\textbf{Herramientas típicas:}

\begin{itemize}
    \item OWASP ZAP.
    \item Burp Suite.
    \item Pruebas basadas en OWASP Top 10.
\end{itemize}

% --------------------------------------------------

\subsection{Paso 2: Red sin segmentación → acceso desde soporte a bases de datos de tarjetas}

Se utilizaría una auditoría de \textbf{Seguridad de la Información y Ciberseguridad} centrada en \textbf{Network Security Assessment}.

\textbf{Habría detectado:}

\begin{itemize}
    \item Falta de segmentación mediante VLAN.
    \item Ausencia de arquitectura Zero Trust.
    \item Movimiento lateral libre entre redes críticas.
\end{itemize}

\textbf{CAATs utilizados:}

\begin{itemize}
    \item Nmap → mapeo de red.
    \item Snort / IDS-IPS → detección de intrusiones.
    \item MITRE ATT\&CK → modelado de técnicas de ataque.
\end{itemize}

% --------------------------------------------------

\subsection{Paso 3: Credenciales en texto plano}

\textbf{Escenario:} Script de backup con contraseña de administrador expuesta.

Se utilizaría una auditoría de \textbf{Aplicaciones y Desarrollo (SDLC)} mediante \textbf{Análisis de Código (DevSecOps / Software Security)}.

\textbf{Habría detectado:}

\begin{itemize}
    \item Credenciales hardcodeadas.
    \item Malas prácticas de programación.
    \item Ausencia de revisiones de código.
\end{itemize}

\textbf{CAATs utilizados:}

\begin{itemize}
    \item SAST (Static Application Security Testing).
    \item Análisis de repositorios de código.
    \item Scripts de auditoría en Python y SQL.
\end{itemize}

% --------------------------------------------------

\subsection{Gestión de Proyectos IT (fallo estructural)}

Este problema constituye la causa raíz organizativa y no forma parte directamente del patrón de ataque.

\textbf{Problema detectado:}

Portal sin actualizar durante dos años.

\textbf{Auditoría aplicable:}

Gestión de Proyectos IT.

\textbf{Habría detectado:}

\begin{itemize}
    \item Falta de mantenimiento del software.
    \item Mala planificación de actualizaciones.
    \item Ausencia de controles de calidad en las entregas.
\end{itemize}

\textbf{Resultado:}

El sistema no habría quedado obsoleto y la superficie de ataque se habría reducido significativamente.

% ==================================================

\section{Normativa}

\subsection{¿Qué auditoría específica acaba de reprobar estrepitosamente la empresa en el paso 4, enfrentándose a multas millonarias?}

\begin{tcolorbox}
\textbf{Auditoría fallida: Auditoría de Cumplimiento Normativo (Compliance)}
\end{tcolorbox}

\textbf{Normativas incumplidas:}

\begin{itemize}
    \item PCI-DSS.
    \item RGPD / LOPDGDD.
    \item ISO 27001.
\end{itemize}

\textbf{Consecuencias:}

\begin{itemize}
    \item Multas económicas elevadas.
    \item Investigación gubernamental.
    \item Pérdida de licencias para el procesamiento de pagos.
\end{itemize}

% ==================================================

\section{Reacciones}

\subsection{Respuesta al Director General (CEO)}

El CEO propone contratar urgentemente una Auditoría de Seguridad Física para evitar que el incidente vuelva a ocurrir.

Como Auditor Jefe, respondería que la auditoría de seguridad física no es prioritaria en este caso, ya que el origen del ataque es exclusivamente lógico y se encuentra en fallos relacionados con el SDLC, la segmentación de red y la gestión de credenciales.

Por ejemplo, el fallo de SDLC se observa cuando una aplicación web se publica con vulnerabilidades como XSS descritas en OWASP, equivalente a dejar una puerta abierta en una aplicación accesible desde Internet.

La falta de segmentación de red permite que, una vez comprometido el sistema de soporte, el atacante alcance bases de datos críticas. Es comparable a un edificio donde todas las salas pueden abrirse sin restricciones. Una arquitectura Zero Trust evitaría este comportamiento.

La mala gestión de credenciales se produce al almacenar contraseñas en texto plano, situación equivalente a dejar una llave maestra visible y accesible para cualquier persona.

Por tanto, la respuesta adecuada consiste en realizar una auditoría integral de ciberseguridad basada en:

\begin{itemize}
    \item OWASP.
    \item Zero Trust.
    \item DevSecOps.
    \item PCI-DSS.
\end{itemize}

La seguridad física únicamente tendría un papel complementario dentro de una estrategia global de protección.

\vfill

\begin{center}
\rule{0.4\textwidth}{0.4pt}

\vspace{0.3cm}

\textbf{Fin del informe}
\end{center}

\end{document}